\documentclass[12pt,a4paper]{article}
\usepackage{graphicx}
\usepackage{titlesec}
\usepackage{setspace}
\usepackage{geometry}
\geometry{margin=1in}

\usepackage{graphicx}
\usepackage{subcaption}
\usepackage{float}
\usepackage{listings}

% Formatting
\setstretch{1.15}
\titleformat{\section}{\large\bfseries}{\thesection.}{1em}{}
\titleformat{\subsection}{\normalsize\bfseries}{\thesubsection.}{1em}{}
\usepackage[colorlinks=true, linkcolor=blue, citecolor=blue, urlcolor=blue]{hyperref}

\lstset{
    basicstyle=\ttfamily\small,
    numbers=left,
    numberstyle=\tiny,
    frame=single,
    breaklines=true,
    showstringspaces=false,
    tabsize=4
}


%----------------------------------------
\begin{document}

% --------- First Page (Cover Page) ---------
\begin{titlepage}
    \centering
    \includegraphics[width=3cm]{Logo.jpeg}\par\vspace{1cm}
    {\LARGE \textbf{Department of Electrical and Information Engineering}}\\[0.4cm]
    {\Large \textbf{University of Ruhuna}}\\[5.5cm]
    {\Large \textbf{Assessment 01}}\\[0.2cm]
    {\large \textbf{Image Processing and Computer Vision}}\\[0.2cm]
    {\large \textbf{EE7204 / EC7205}}\\[6cm]
    
   % Centered student and project info
    {\large \textbf{M.R.M.ASHFAQ}}\\[0.4cm]
    {\large \textbf{EF/2021/4417}}\\[0.4cm]
    {\large \textbf{Computer Engineering (CE)}}\\[0.4cm]
   
    {\large \textbf{01/03/2026}}\\[2cm]
    
\end{titlepage}


% --------- Table of Contents (Clickable) ---------

\clearpage
\phantomsection
\addcontentsline{toc}{section}{Contents}
\label{contents}
\tableofcontents
\clearpage

% Marks 
\section*{Obtained Marks Distribution}

\begin{table}[H]
\centering
\renewcommand{\arraystretch}{2.5} % increase row height
\begin{tabular}{|c|l|c|}
\hline
\textbf{Q. No} & \textbf{Description} & \textbf{Marks} \\
\hline
1 & Preliminary Question 1 &  \\
2 & Preliminary Question 2 & \\
3 & Preliminary Question 3 &  \\
4 & Preliminary Question 4 &  \\
5 & Preliminary Question 5 &  \\
6 & Preliminary Question 6 &  \\
7 & Preliminary Question 7 &  \\
8 & Preliminary Question 8 &  \\
9 & Preliminary Question 9 &  \\
10 & Preliminary Question 10 &  \\
\hline
11 & Practical Part (1 Question) &  \\
\hline
\multicolumn{2}{|c|}{\textbf{Total Marks}} & \textbf{} \\
\hline
\end{tabular}
\caption{Marks allocation for preliminary and practical parts}
\label{tab:marks_per_question}
\end{table}

\newpage{}



% --------- Sections ---------

\section{Preliminary work}

\subsection{Question 1}
\subsubsection*{Snapshot of the code}
\lstinputlisting[language=Python, caption={Code snapshot for Question 1 (EG\_2021\_4417\_Q1.py)}, label={lst:q1_code}]{../Codes/EG_2021_4417_Q1.py}

\subsubsection*{Results / outputs}
\begin{figure}[H]
    \centering
    \begin{subfigure}[b]{0.30\textwidth}
        \centering
        \includegraphics[width=\textwidth]{../Images/Image_1.jpg}
        \caption{Original}
        \label{fig:q1_original}
    \end{subfigure}
    \hfill
    \begin{subfigure}[b]{0.30\textwidth}
        \centering
        \includegraphics[width=\textwidth]{../Outputs/Q1/Image_1_avg_3x3.png}
        \caption{Average $3\times 3$}
        \label{fig:q1_avg_3x3}
    \end{subfigure}
    \hfill
    \begin{subfigure}[b]{0.30\textwidth}
        \centering
        \includegraphics[width=\textwidth]{../Outputs/Q1/Image_1_avg_5x5.png}
        \caption{Average $5\times 5$}
        \label{fig:q1_avg_5x5}
    \end{subfigure}

    \vspace{6pt}

    \begin{subfigure}[b]{0.30\textwidth}
        \centering
        \includegraphics[width=\textwidth]{../Outputs/Q1/Image_1_avg_11x11.png}
        \caption{Average $11\times 11$}
        \label{fig:q1_avg_11x11}
    \end{subfigure}
    \hfill
    \begin{subfigure}[b]{0.30\textwidth}
        \centering
        \includegraphics[width=\textwidth]{../Outputs/Q1/Image_1_avg_15x15.png}
        \caption{Average $15\times 15$}
        \label{fig:q1_avg_15x15}
    \end{subfigure}
    \caption{Average filtering results for different kernel sizes.}
    \label{fig:q1_results}
\end{figure}

\subsubsection*{Discussion / observations}
Average filtering is a linear low-pass operation that reduces high-frequency components (noise and texture) by replacing each pixel with the mean of its neighbourhood.
As the kernel size increases from $3\times 3$ to $15\times 15$, the smoothing effect becomes stronger: noise is reduced more, but edges become progressively blurred and fine details are lost.
In this experiment, $3\times 3$ and $5\times 5$ preserve most structures with mild smoothing, while $11\times 11$ and $15\times 15$ noticeably reduce sharpness and local contrast.

\hyperref[contents]{\small Back to Table of Contents}
\subsection{Question 2}

\subsubsection*{Snapshot of the code}
\lstinputlisting[language=Python, caption={\texttt{EG\_2021\_4417\_Q2.py}}, label={lst:q2_code}]{../Codes/EG_2021_4417_Q2.py}

\subsubsection*{Results / outputs}
\begin{figure}[H]
    \centering
    \begin{subfigure}[b]{0.32\textwidth}
        \centering
        \includegraphics[width=\textwidth]{../Images/Image_2.jpg}
        \caption{Original}
        \label{fig:q2_original}
    \end{subfigure}
    \hfill
    \begin{subfigure}[b]{0.32\textwidth}
        \centering
        \includegraphics[width=\textwidth]{../Outputs/Q2/Image_2_salt_pepper_10pct.png}
        \caption{S\&P $10\%$}
        \label{fig:q2_sp_10}
    \end{subfigure}
    \hfill
    \begin{subfigure}[b]{0.32\textwidth}
        \centering
        \includegraphics[width=\textwidth]{../Outputs/Q2/Image_2_salt_pepper_20pct.png}
        \caption{S\&P $20\%$}
        \label{fig:q2_sp_20}
    \end{subfigure}
    \caption{Salt \& pepper noise added to \texttt{Image\_2.jpg}.}
    \label{fig:q2_noise}
\end{figure}

\begin{figure}[H]
    \centering
    \begin{subfigure}[b]{0.32\textwidth}
        \centering
        \includegraphics[width=\textwidth]{../Outputs/Q2/Image_2_salt_pepper_10pct_median_3x3.png}
        \caption{$10\%$ + Median $3\times 3$}
        \label{fig:q2_10_m3}
    \end{subfigure}
    \hfill
    \begin{subfigure}[b]{0.32\textwidth}
        \centering
        \includegraphics[width=\textwidth]{../Outputs/Q2/Image_2_salt_pepper_10pct_median_5x5.png}
        \caption{$10\%$ + Median $5\times 5$}
        \label{fig:q2_10_m5}
    \end{subfigure}
    \hfill
    \begin{subfigure}[b]{0.32\textwidth}
        \centering
        \includegraphics[width=\textwidth]{../Outputs/Q2/Image_2_salt_pepper_10pct_median_11x11.png}
        \caption{$10\%$ + Median $11\times 11$}
        \label{fig:q2_10_m11}
    \end{subfigure}

    \vspace{6pt}

    \begin{subfigure}[b]{0.32\textwidth}
        \centering
        \includegraphics[width=\textwidth]{../Outputs/Q2/Image_2_salt_pepper_20pct_median_3x3.png}
        \caption{$20\%$ + Median $3\times 3$}
        \label{fig:q2_20_m3}
    \end{subfigure}
    \hfill
    \begin{subfigure}[b]{0.32\textwidth}
        \centering
        \includegraphics[width=\textwidth]{../Outputs/Q2/Image_2_salt_pepper_20pct_median_5x5.png}
        \caption{$20\%$ + Median $5\times 5$}
        \label{fig:q2_20_m5}
    \end{subfigure}
    \hfill
    \begin{subfigure}[b]{0.32\textwidth}
        \centering
        \includegraphics[width=\textwidth]{../Outputs/Q2/Image_2_salt_pepper_20pct_median_11x11.png}
        \caption{$20\%$ + Median $11\times 11$}
        \label{fig:q2_20_m11}
    \end{subfigure}
    \caption{Median filtering results for different noise levels and kernel sizes.}
    \label{fig:q2_median_results}
\end{figure}

\subsubsection*{Discussion / observations}
Salt \& pepper noise is impulse noise that appears as random black/white pixels; increasing the corruption rate from $10\%$ to $20\%$ significantly increases the speckle density and obscures fine details.
Median filtering is effective for impulse noise because it removes outlier pixels while preserving edges better than linear averaging.
However, as the kernel size increases (e.g., $11\times 11$), denoising improves but edges and thin structures can become distorted and textures become over-smoothed.

\hyperref[contents]{\small Back to Table of Contents}
\subsection{Question 3}

\hyperref[contents]{\small Back to Table of Contents}
\subsection{Question 4}

\hyperref[contents]{\small Back to Table of Contents}
\subsection{Question 5}

\hyperref[contents]{\small Back to Table of Contents}
\subsection{Question 6}

\hyperref[contents]{\small Back to Table of Contents}
\subsection{Question 7}

\hyperref[contents]{\small Back to Table of Contents}
\subsection{Question 8}

\hyperref[contents]{\small Back to Table of Contents}
\subsection{Question 9}

\hyperref[contents]{\small Back to Table of Contents}
\subsection{Question 10}


\hyperref[contents]{\small Back to Table of Contents}
\newpage


\section{Practical Work}
\subsection{Introduction}
(Maximum word limit: 50 words)


\subsection{Flow Diagram of the image processing pipeline }
\begin{figure}[H]
    \centering
    \includegraphics[width=0.70\textwidth]{01_Original.png}
    \caption{Flow Diagram}
    \label{fig:flow_diagram}
\end{figure}



\hyperref[contents]{\small Back to Table of Contents}
\newpage
\section{Introduction and stepwise justification of the selected image processing pipeline.}
(Maximum word limit: 100 words)



\hyperref[contents]{\small Back to Table of Contents}
\newpage
\section{Expected Results for 10 Images}

\begin{figure}[H]
\centering
\renewcommand{\arraystretch}{1.2}
\begin{tabular}{ccc}

\textbf{Original} & \textbf{Result} & \textbf{Ground Truth} \\[6pt]

\includegraphics[width=0.28\textwidth]{01_Original.png} &
\includegraphics[width=0.28\textwidth]{01_Original.png} &
\includegraphics[width=0.28\textwidth]{01_Original.png} \\[10pt]

% -------- Uncomment when images are available --------
% \includegraphics[width=0.28\textwidth]{img2_original.png} &
% \includegraphics[width=0.28\textwidth]{img2_result.png} &
% \includegraphics[width=0.28\textwidth]{img2_gt.png} \\[10pt]
%
% \includegraphics[width=0.28\textwidth]{img3_original.png} &
% \includegraphics[width=0.28\textwidth]{img3_result.png} &
% \includegraphics[width=0.28\textwidth]{img3_gt.png} \\[10pt]
%
% ...
%
% \includegraphics[width=0.28\textwidth]{img10_original.png} &
% \includegraphics[width=0.28\textwidth]{img10_result.png} &
% \includegraphics[width=0.28\textwidth]{img10_gt.png}

\end{tabular}
\caption{Comparison of original images, obtained results, and ground truth for 10 test cases}
\label{fig:results_vs_gt}
\end{figure}





\hyperref[contents]{\small Back to Table of Contents}
\newpage
\section{Results}
Average results obtained from the 50-image validation set are presented in Table 1.


\hyperref[contents]{\small Back to Table of Contents}
\newpage
\section{Discussion}
{Maximum Words limits : 100 Words}



\bibliographystyle{ieeetr} 
\bibliography{references}
% Add the references if any.


\end{document}
